\section{Limitations and Ethics}
\label{sec:limitations}

The present evidence should be interpreted within four boundaries. First, the synthetic web-intelligence text is LLM-generated with an \texttt{as\_of\_date} constraint, so residual parametric leakage cannot be fully ruled out; we therefore use the synthetic split mainly for controlled development and mechanism probing. Second, the real-news corpus is the paper's primary source of external-validity evidence, but it spans only December 2025--March 2026 and the ICAIF main comparison uses two seeds and five broad test blocks, so confidence intervals remain wide and the results should be read as support for problem validity and model plausibility rather than as benchmark-complete or cross-regime evidence. Third, the downstream trading evaluation now includes a 5 bps per-trade cost and fee sensitivity, but it still omits several live-market frictions such as order-book depth, latency, funding, liquidation constraints, and adaptive position sizing; Sharpe should therefore be treated as an offline utility proxy rather than a live-trading projection. Fourth, the task is defined on news-available bars with valid aligned web context rather than on every market bar, making the current conclusions specific to an event-conditioned asynchronous setting. Any deployment in automated trading would require substantial additional validation, risk controls, regulatory compliance, and human oversight.
